% Lösungen Einheit 4
\einheitenkopf{Einheit 4 von 5 · Graph und Ableitungsgraph}
\erg{31}{a) $x = -2$, $x = 0$, $x = 2$ \quad b) $-$, $+$, $-$, $+$ \quad c) $x = 0$, $x = 2$ \quad d) $f_2'(-1) > 0$, $f_2'(1) < 0$, $f_2'(3) > 0$ \quad e) $W(1 \mid 1)$; der Graph fällt dort (zwischen Hochpunkt und Tiefpunkt), also ist die Steigung negativ; $f_2'(1) = -1{,}5$ \quad f) ein Punkt zwischen $x = 1$ und $x = 2$, z.\,B. bei $x = 1{,}5$ (Graph fällt, ist linksgekrümmt und liegt über der $x$-Achse)}
\erg{32}{a) $f'(x) = x^2 - 2x$; Werte $3$; $0$; $-1$; $0$; $3$ \quad b) Parabel durch $(-1 \mid 3)$, $(0 \mid 0)$, $(1 \mid -1)$, $(2 \mid 0)$, $(3 \mid 3)$ \quad c) Im Hochpunkt $(0 \mid 0)$ und im Tiefpunkt $(2 \mid -\frac{4}{3})$ ist die Tangente waagerecht, also $f' = 0$; dazwischen fällt $f$, dort ist $f'$ negativ.}
\erg{33}{a) und b) mögliche Skizzen:}

\begin{ksys}[xmin=-4,xmax=4,ymin=-4,ymax=3,klein]
\funktionab{5/32*\x^3-15/8*\x-0.5}{f}{-4}{4}
\hochpunkt{-2}{2}{H}
\tiefpunkt{2}{-3}{T}
\end{ksys}\ksysabstand\begin{ksys}[xmin=-4,xmax=4,ymin=-3,ymax=3,klein]
\funktionab{15/32*\x^2-15/8}{f'}{-4}{4}
\end{ksys}

\noindent $f'$: Parabel mit den Nullstellen $-2$ und $2$ und dem Tiefpunkt bei $x = 0$ \quad c) Zwei Tiefpunkte und ein Hochpunkt: $f'$ hat mindestens drei Nullstellen mit Vorzeichenwechsel, also mindestens Grad 3; $f$ hat damit mindestens Grad 4.

\erg{34}{a) $-1 < x < 3$ \quad b) $x = -1$ und $x = 3$ \quad c) bei $x = 1$, der Abstand ist $d(1) = 2$}
\erg{35}{a) Wendepunkt $W(1 \mid 2)$, $f'(1) = 1{,}5$; Normale: Steigung $-\frac{1}{1{,}5} = -\frac{2}{3}$ – Aussage (1) stimmt. $f'$ ist eine nach unten geöffnete Parabel mit Scheitel $(1 \mid 1{,}5)$, Wertebereich $]{-\infty};\,1{,}5]$ – Aussage (2) ist falsch, Werte über $1{,}5$ nimmt $f'$ nicht an. \quad b) falsch: Die Wendetangente $y = 1{,}5x + 0{,}5$ hat mit dem Graphen nur den Wendepunkt gemeinsam, denn $f(x) - (1{,}5x + 0{,}5) = -0{,}5(x - 1)^3$; alle anderen Tangenten haben zwei gemeinsame Punkte.}
\erg{36}{a) Ben liest den Graphen von $f'$ wie den von $f$. Der Tiefpunkt von $f'$ bei $x = 1$ ist die Wendestelle von $f$ (dort fällt $f$ am stärksten). \quad b) $f'(-1) = 0$ mit Wechsel von $+$ nach $-$: Hochstelle bei $x = -1$; $f'(3) = 0$ mit Wechsel von $-$ nach $+$: Tiefstelle bei $x = 3$}
\erg{37}{a) Nach der Potenzregel wird aus $a\,x^n$ der Summand $n\,a\,x^{n-1}$; die höchste Potenz sinkt um eins. \quad b) Im Wendepunkt wechselt $f''$ das Vorzeichen; $f''$ ist die Ableitung von $f'$, also wechselt $f'$ dort von Steigen zu Fallen oder umgekehrt: $f'$ hat dort einen Hoch- oder Tiefpunkt. \quad c) Angenommen, die Graphen schneiden sich in $[0;\,4]$ nicht. Da beide Graphen ohne Sprung verlaufen, liegt dann einer überall über dem anderen, und seine Fläche ist größer. Das widerspricht den gleich großen Flächen; also gibt es mindestens einen Schnittpunkt.}
\erg{38}{a) von 2 Uhr bis 10 Uhr ($2 < t < 10$, dort $w'(t) > 0$) \quad b) Tiefpunkt bei $t = 2$ ($w'$ wechselt von $-$ nach $+$), Hochpunkt bei $t = 10$ ($w'$ wechselt von $+$ nach $-$) \quad c) bei $t = 6$, also um 6 Uhr, mit 4 cm pro Stunde}
